\subsection*{11. The Founder's Unpaid Work Is a Subsidy}

A founder keeping a full-time job may sensibly prioritize cash survival, but economic analysis must not confuse an unbilled hour with a free hour. Founder work can be priced in at least 3 ways: the cost of replacing it with a competent person, the founder's actual alternative earning opportunity, or the cost of losing capacity for other essential work. These prices answer different questions.

Build 2 parallel ledgers:

\begin{table*}[t]
\centering
\small
\caption{Cash Pilot Ledger Versus Normalized Repeatability Ledger}
\begin{tabularx}{\textwidth}{p{0.20\textwidth}p{0.36\textwidth}X}
\toprule
\textbf{Cost Component} & \textbf{Cash Pilot Ledger} & \textbf{Normalized Repeatability Ledger} \\
\midrule
Fee collected & Count actual receipts and obligations & Count sustainable realized revenue \\
Teaching by founder & ₹0 cash salary if unpaid & Competent replacement cost or appropriate opportunity cost \\
Curriculum development & Cash outlay now & Amortized or expensed using a defensible planning horizon \\
Mentor and support & Actual cash paid & Full staffing needed at promised standard \\
Acquisition & Ad spend and paid tools & Founder sales hours plus tools and staff \\
Quality failures & Actual refunds and rework & Expected future recovery and replacement costs \\
\bottomrule
\end{tabularx}
\end{table*}

The goal is not to make a pilot appear unprofitable through arbitrary imaginary salaries. It is to avoid building a ``scalable'' model that depends permanently on founder labor that cannot be replicated.

\subsection*{12. Direct Costs and Allocations Can Tell Contradictory Stories}

A college partnership and a D2C cohort might share curriculum, an LMS, instructor training, and the same faculty. Allocate every shared rupee equally and 1 service may look terrible; allocate none and both may look unusually attractive. Neither tells the whole story.

For a decision to accept an additional contract, ask: What \textbf{incremental costs} will it impose? What scarce capacity will it displace? What customer obligations will persist afterward? For a decision to keep a business line long term, ask: Does it contribute enough to sustain \textbf{shared infrastructure} and eventual replacement of people and assets? For pricing, ask which costs grow if the segment succeeds.

An activity-based costing lens: mapping major cost drivers to activities such as admissions calls, project review, deployment, and escalation: is useful. It need not become a complicated software project. A simple estimate of \textbf{minutes per consequential activity} is often more revealing than a polished full-cost spreadsheet.

\subsection*{13. Cost per Outcome and the Selection Problem}

Suppose a selective program has more successful graduates than an open-admission 1. Its lower cost per success may reflect a stronger entry population rather than superior instructional efficiency. Similarly, an enterprise intervention may look cheap because the client already possessed a competent technical team and usable data.

Keep the customer starting condition in the ledger. Compare like with like where feasible: cost per successfully served \textbf{relevant learner segment}, time to mastery for comparable prerequisites, feedback minutes per corrected misconception, or cost per verified task improvement. This is an operational comparison, not proof of causal instructional superiority.

Do not ``improve'' efficiency by admitting only learners who barely need the service and then crediting the institution for their outcomes. Selection has an honest role in creating a coherent cohort, but it must not masquerade as learning productivity.

\subsection*{14. The Cash-Flow Paradox of Prepaid Services}

Prepayment can finance development; it also creates an obligation to deliver later. A program that collects ₹6 lakh before classes begin may appear flush with cash while owing months of instruction, support, and potential refunds. A college contract booked for ₹20 lakh can create a different problem: significant instructor and travel costs may be due before the university pays its invoice.

A founder should chart \textbf{cash in, cash out, service milestones, and refund exposure by week or month}. That is an operations document as much as a financial document. When a business collects early, preserve a reserve corresponding to remaining obligations; do not treat every receipt as spendable surplus.

This principle becomes vital when educational businesses expand across cohorts: cash from the next intake must not become the only way to finish the previous intake's promised teaching.
